\section{Product realizations and their measures}
\label{sec:trace-product-measures}

For a sum of potentials pulled back from the two factors of a product with a curve, dimension reasons force every term of the Monge--Amp\`ere expansion to contain exactly one factor from the curve. The resulting product formula computes the measure of the pulled-back Green representative and solves equations with separated densities. Adding a potential from the other factor with trivial multiplier ideals can keep every trace coefficient while changing finite mass to infinite mass.

In this section, $C$ is a compact connected curve with a smooth positive form $\eta$ of integral one, and $(Z,\omega_Z)$ is a compact connected Hermitian manifold of dimension $k\geq1$. On $Y=C\times Z$, let $q:Y\to C$ and $r:Y\to Z$ be the projections. Choose $A>0$ and set
\[
 \Omega=q^*(A\eta)+r^*\omega_Z.
\]
For positive Borel measures $\mu$ on $C$ and $\nu$ on $Z$, $\mu\otimes\nu$ denotes their product measure. The extended measures below are $\sigma$-finite by \cref{thm:arbitrary-extended-ma}, so this notation is meaningful also when a total mass is infinite.

\begin{proposition}\label{prop:curve-product-extended-ma}\index{Trace!product measure}
For $w\in\PSH(C,A\eta)$ and $h\in\PSH(Z,\omega_Z)$, the potential $a=w\circ q+h\circ r$ is $\Omega$-plurisubharmonic and
\begin{equation}\label{eq:curve-product-extended-ma}
 \mathrm{MA}^{\mathrm{ext}}_\Omega(a)
 =(k+1)\,
       \mathrm{MA}^{\mathrm{ext}}_{A\eta}(w)
                \otimes\mathrm{MA}^{\mathrm{ext}}_{\omega_Z}(h).
\end{equation}
\end{proposition}
\begin{proof}
The curvature current of $a$ is the sum of the pullbacks of the two positive currents $A\eta+\ddc w$ and $\omega_Z+\ddc h$, which proves positivity.

Suppose first that both potentials are bounded. Expanding the bounded product by multilinearity, only the terms with one factor from $C$ and $k$ factors from $Z$ survive: two factors from $C$ vanish for dimension reasons, and so do $k+1$ factors from $Z$. The binomial coefficient is $k+1$, which gives the formula for bounded potentials once the exterior product $q^*(A\eta+\ddc w)\wedge r^*(\omega_Z+\ddc h)^k$ is identified with the product measure $(A\eta+\ddc w)\otimes(\omega_Z+\ddc h)^k$. To justify this identification at the present regularity, add local plurisubharmonic corrections on each factor and regularize by decreasing convolution on smaller product charts. For the smooth approximants the identity is the exterior product of forms. Local Bedford--Taylor convergence, together with the finite expansion in the smooth background, gives weak convergence of the factor measures, with uniformly bounded variation on compact subcharts. The product measures then converge weakly: one first tests against products of continuous functions on the two factors and then uses uniform approximation by finite sums of such products. This proves the identity in the limit.

For general potentials, subtract constants so that $w,h\leq0$, and set $w_j=\max(w,-j)$ and $h_j=\max(h,-j)$. On
\[
 R_j=\{w>-j\}\times\{h>-j\}
\]
the sum agrees with the bounded potential $w_j\circ q+h_j\circ r$. This set is plurifine open and Borel, and it lies in $\{a>-2j\}$. Locality of the products of bounded truncations therefore identifies the restriction to $R_j$ of the extended measure of $a$ with that of the bounded sum. By the bounded case, this is $k+1$ times the product of the restrictions of the truncated measures, which are exactly the restrictions of the two extended measures of the factors.

The sets $R_j$ increase to the product of the complements of the two pole sets, and both sides of \eqref{eq:curve-product-extended-ma} give zero mass to the remaining set: for the left-hand side it is the pole set of $a$, and for the right-hand side this follows because each factor measure charges no pluripolar set and the product measures are $\sigma$-finite. Monotone convergence of positive measures proves the formula, even when the total mass is infinite.
\end{proof}

Suppose now that $D$ is a finite atomic measure on $C$, for instance the trace measure of \cref{sec:curve-fibration-traces}, put $\delta=D(C)$, and let $w$ be the Green potential with $\ddc w=D-\delta\eta$. For $A>\delta$, the potential $v=w\circ q$ satisfies
\begin{equation}\label{eq:product-trace-smooth-measure}
 \mathrm{MA}^{\mathrm{ext}}_\Omega(v)
       =(k+1)(A-\delta)\,\eta\otimes\omega_Z^k.
\end{equation}
Thus its measure is smooth and positive, although $v$ may have infinitely many singular fibers. This identity concerns the extended measure, whose definition discards the divisorial part of the curvature current.

The product formula also solves equations with separated densities in this ordinary type. Fix a smooth positive volume form $\mathrm{d}V_Z$ on $Z$ and $p>1$, and let $0\leq f\in L^p(C,\eta)$ and $0\leq g\in L^p(Z,\mathrm{d}V_Z)$ have positive integrals. By \cref{thm:curve-atomic-lp-equation}, applied on $C$ in the background $A\eta$, there is $w_f\sim w$ with
\[
 \mathrm{MA}^{\mathrm{ext}}_{A\eta}(w_f)=c_C f\eta,
 \qquad c_C=\frac{A-\delta}{\int_Cf\eta}.
\]
\Needspace{6\baselineskip}
By \cref{thm:kn-continuous}, there are a bounded $h_g\in\PSH(Z,\omega_Z)$ and $c_Z>0$ with
\[
 \mathrm{MA}_{\omega_Z}(h_g)=c_Zg\,\mathrm{d}V_Z.
\]
The sum $a=w_f\circ q+h_g\circ r$ is ordinarily equivalent to $v$, and
\[
 \mathrm{MA}^{\mathrm{ext}}_\Omega(a)
       =(k+1)c_Cc_Z\,(f\eta)\otimes(g\,\mathrm{d}V_Z).
\]
The product density belongs to $L^p$ by Fubini's theorem. This is an explicit construction for separated densities; existence for general $L^p$ densities, in a sufficiently large background, follows from \cref{cor:trace-envelope-prescribed-equation}.

\begin{proposition}\label{prop:exact-trace-infinite-representatives}
Suppose that $Z$ does not have the bounded mass property, that is, $v_+(\omega_Z)=+\infty$ in the notation of \cref{def:bmp}, and that $A>\delta$. Then there is $a\in\PSH(Y,\Omega)$ such that
\[
 a\sim_{\mathcal I}v,\qquad
 \int_Y\mathrm{MA}^{\mathrm{ext}}_\Omega(v)<+\infty,
 \qquad
 \int_Y\mathrm{MA}^{\mathrm{ext}}_\Omega(a)=+\infty.
\]
Both potentials have the same divisorial Lelong numbers on every modification; when $D$ is the trace measure of \cref{sec:curve-fibration-traces}, these are the trace data computed in \cref{thm:curve-fibration-exact-trace}. The product manifold $Y$ does not have the bounded mass property.
\end{proposition}
\begin{proof}
Choose the nonpositive potential $H$ of \cref{thm:infinite-mass-witness} on $Z$; it has infinite extended mass, and $\mathrm{e}^{-tH}$ is integrable for every $t>0$. Put $a=w\circ q+H\circ r$. The product formula and $A-\delta>0$ show that the mass of $a$ is infinite, while \eqref{eq:product-trace-smooth-measure} gives the finite positive mass of $v$.

We prove equality of the multiplier ideals locally, for every $t>0$. Since $H\leq0$, integrability against $\mathrm{e}^{-ta}$ implies integrability against $\mathrm{e}^{-tv}$, so $\mathcal I(ta)\subseteq\mathcal I(tv)$. Conversely, \cref{lem:submersion-multiplier-pullback} gives $\mathcal I(tv)=\mathcal I(tw)\mathcal O_Y$. Local generators pulled back from $C$ are square integrable against $\mathrm{e}^{-tw}$, so by Fubini's theorem they are also square integrable against $\mathrm{e}^{-tw}\mathrm{e}^{-tH}=\mathrm{e}^{-ta}$; multiplication by bounded holomorphic functions and finite sums preserve this local integrability. Thus $\mathcal I(tv)\subseteq\mathcal I(ta)$, which proves equality at every positive weight. \Cref{thm:local-i-inputs} then gives equality of all divisorial Lelong numbers after modification.

Finally, since $a$ has infinite mass, \cref{prop:bmp-finite} shows that $Y$ does not have the bounded mass property.
\end{proof}

For a concrete ambient realization, take a compact complex manifold $M$ containing the curve $C$, a potential $U\in\QPSH(M)$ with $\nu(U,C)=0$, and $X=M\times Z$. Then $Y=C\times Z$ and $u=U\circ\operatorname{pr}_M$, where $\operatorname{pr}_M:X\to M$ is the projection, satisfy the hypotheses of \cref{thm:curve-fibration-exact-trace} with $p=\operatorname{pr}_M$. The factor $Z$ can be the Hopf threefold of \cref{thm:hopf-unbounded}. Small logarithmic poles of $U$ at finitely many points of $C$ already give nonzero trace data, and the construction on the curve also allows any countable trace data arising from the multiplier ideals of $U$.

The mass formula of \cref{sec:resolved-trace-masses} uses more information: the restriction itself in a closed background, with a finite divisorial part and a bounded remainder on a K\"ahler model. Its residual class is nef and computes the mass of that restriction. Divisorial trace data alone does not provide such a remainder, as the preceding example shows.
